\documentclass[12pt]{article}
\usepackage[spanish]{babel}
\usepackage[utf8]{inputenc}
\usepackage[T1]{fontenc}
\usepackage{amsmath, amssymb}
\usepackage{geometry}
\usepackage{enumitem}
\usepackage{needspace}
\usepackage{tikz}
\usetikzlibrary{babel}

\geometry{letterpaper, margin=2cm}
\setlength{\parindent}{0pt}
\setlength{\parskip}{6pt}

\newcommand{\opcion}[2]{\item[#1)] $#2$}
\newcommand{\puntografico}{0.35}
\newcommand{\puntograficoafirmacioncuatro}{0.09}
\newcommand{\puntograficoitemunoafirmacioncuatro}{0.14}
\newcommand{\grafsenal}[4]{%
\begin{tikzpicture}[scale=0.09]
  \path[use as bounding box] (-15,-22) rectangle (32,24);
  \draw[<->, thick] (-13,0) -- (30,0);
  \draw[<->, thick] (0,-20) -- (0,21);
  \node[anchor=west, inner sep=1pt] at (30.05,0) {\scriptsize $x$ (este)};
  \node[anchor=south, inner sep=1pt] at (0,21.55) {\scriptsize $y$ (norte)};
  \draw[thick] (#1,0) circle (#2);
  \fill (0,0) circle (\puntografico) node[below left] {\scriptsize $M$};
  \fill (#1,0) circle (\puntografico) node[above] {\scriptsize $A$};
  \fill (#1+#2,0) circle (\puntografico);
  \node[anchor=north, xshift=-2pt] at (#1,0) {\scriptsize $#4$};
  \node[anchor=north, xshift=3pt] at (#1+#2,0) {\scriptsize $#3$};
\end{tikzpicture}}
\newcommand{\graftrasbase}[7]{%
\begin{tikzpicture}[scale=0.16]
  \path[use as bounding box] (-11,-9) rectangle (31,24);
  \draw[<->, thick] (-8,0) -- (28,0);
  \draw[<->, thick] (0,-7) -- (0,21);
  \node[anchor=west, inner sep=1pt] at (28.05,0) {$x$};
  \node[anchor=south, inner sep=1pt] at (0,21.05) {$y$};
  \draw[thick] (#4,#5) circle (#3);
  \draw[dashed] (#4,0) -- (#4,#5);
  \draw[dashed] (0,#5) -- (#4,#5);
  \fill (#4,#5) circle (\puntografico) node[anchor=south west, inner sep=2pt] {$A'$};
  #7
  \node[left] at (0,#5) {$#5$};
\end{tikzpicture}}
\newcommand{\graftras}[6]{%
\graftrasbase{#1}{#2}{#3}{#4}{#5}{#6}{\node[anchor=north] at (#4,0) {$#4$};}}
\newcommand{\graftrasetiquetabaja}[6]{%
\graftrasbase{#1}{#2}{#3}{#4}{#5}{#6}{\node[anchor=north, inner sep=1pt] at (#4,#5-0.45) {$#4$};}}
\newcommand{\tablarel}[2]{%
\begin{tabular}{|c|*{5}{c|}}
\hline
#1 \\ \hline
#2 \\ \hline
\end{tabular}}
\newcommand{\puntoabierto}[2]{\draw[fill=white, thick] (#1,#2) circle (0.15);}
\newcommand{\puntocerrado}[2]{\fill (#1,#2) circle (0.15);}

\begin{document}

\begin{center}
  {\Large \textbf{Primera Pr\'actica de la Prueba Nacional}}\\
  {\Large \textbf{Estandarizada Sumativa de Matem\'atica 2026}}\\[6pt]
  {\scriptsize\scshape Elaborado por: Alberto Jim\'enez Ruiz}\\[-1pt]
  {\scriptsize Asesor Regional de Matem\'atica}\\[-1pt]
  {\scriptsize DRE-Cartago 2026 \textbar{} DAP}\\[5pt]
  Nombre: \rule{8cm}{0.4pt}\hfill Fecha: \rule{3cm}{0.4pt}
\end{center}

\section*{Instrucciones}
Seleccione la opci\'on que corresponde a la representaci\'on algebraica o gr\'afica descrita en cada situaci\'on. Todas las unidades est\'an dadas en metros, salvo que se indique otra unidad.

\begin{enumerate}[label=\textbf{\arabic*.}, itemsep=1.1cm]

\Needspace{0.92\textheight}
\item
Una antena ($A$) que emite una se\~nal para radio se ubica a $8$ km al este del mercado ($M$) de un poblado, el cual se considera como origen. Si el alcance m\'aximo de la se\~nal que emite la antena es de $10$ km a su alrededor, entonces, la representaci\'on gr\'afica del alcance m\'aximo de la se\~nal corresponde a

\vspace{1\baselineskip}

\noindent\makebox[0pt][l]{\hspace*{1.25cm}%
\setlength{\tabcolsep}{0pt}
\renewcommand{\arraystretch}{1}
\begin{tabular}{@{}r@{\hspace{0.7cm}}c@{}}
\raisebox{1.95cm}{\textbf{A)}} & \grafsenal{8}{5}{13}{8} \\[0.7cm]
\raisebox{1.95cm}{\textbf{B)}} & \grafsenal{8}{10}{18}{8} \\[0.7cm]
\raisebox{1.95cm}{\textbf{C)}} & \grafsenal{8}{20}{28}{8} \\[0.7cm]
\raisebox{1.95cm}{\textbf{D)}} & \grafsenal{10}{8}{18}{10}
\end{tabular}
}


\Needspace{0.90\textheight}
\item
Un dispositivo mec\'anico se coloca en el jard\'in de una casa y lanza un chorro continuo de agua, cuyo alcance m\'aximo es $6$ m alrededor de ese dispositivo.

La siguiente representaci\'on gr\'afica, cuyas unidades est\'an en metros, muestra las ubicaciones, $W$ de esa casa, $D$ del dispositivo y $C$ de la circunferencia correspondiente al alcance m\'aximo del chorro de agua lanzado por ese dispositivo:

\begin{center}
\begin{tikzpicture}[scale=0.15]
  \path[use as bounding box] (-4,-4) rectangle (35,36);
  \draw[<->, thick] (-2,0) -- (32,0);
  \draw[<->, thick] (0,-2) -- (0,34);
  \node[anchor=west, inner sep=1pt] at (32.1,0) {$x$ (este)};
  \node[anchor=south, inner sep=1pt] at (0,34.3) {$y$ (norte)};
  \draw[dashed] (20,0) -- (20,24);
  \draw[dashed] (26,0) -- (26,24);
  \draw[dashed] (0,24) -- (26,24);
  \draw[thick] (20,24) circle (6);
  \fill (0,0) circle (\puntografico) node[below left] {$W$};
  \fill (20,24) circle (\puntografico) node[above] {$D$};
  \fill (26,24) circle (\puntografico) node[above right] {$C$};
  \node[anchor=north] at (20,0) {$20$};
  \node[anchor=north] at (26,0) {$26$};
  \node[left] at (0,24) {$24$};
\end{tikzpicture}
\end{center}

De acuerdo con la informaci\'on anterior, \textquestiondown Cu\'al es la representaci\'on algebraica, cuyas unidades est\'an en metros, de la circunferencia $C$ correspondiente al alcance m\'aximo del chorro de agua lanzado por ese dispositivo?

\begin{enumerate}[label=\Alph*)]
  \opcion{A}{(x-20)^2+(y-24)^2=36}
  \opcion{B}{(x+20)^2+(y+24)^2=36}
  \opcion{C}{(x-20)^2+(y-24)^2=6}
  \opcion{D}{(x-26)^2+(y-24)^2=36}
\end{enumerate}

\Needspace{0.88\textheight}
\item
Una antena de telecomunicaciones ($T$) emite una se\~nal que permite a los tel\'efonos celulares, ubicados a $13$ km o menos alrededor de esa antena, realizar y recibir llamadas. En una representaci\'on gr\'afica, la ubicaci\'on de la antena corresponde al punto $T(5,12)$.

Durante una revisi\'on t\'ecnica se registraron tres tel\'efonos celulares ubicados en los puntos $R(10,24)$, $S(18,12)$ y $U(17,18)$.

La siguiente representaci\'on gr\'afica, en la que las unidades est\'an en kil\'ometros, muestra la ubicaci\'on de la antena y la circunferencia correspondiente al alcance m\'aximo de la se\~nal.

\begin{center}
\begin{tikzpicture}[scale=0.16]
  \path[use as bounding box] (-6,-5) rectangle (27,31);
  \draw[<->, thick] (-4,0) -- (25,0);
  \draw[<->, thick] (0,-3) -- (0,29);
  \node[anchor=west, inner sep=1pt] at (25.05,0) {$x$ (este)};
  \node[anchor=south, inner sep=1pt] at (0,29.55) {$y$ (norte)};
  \draw[thick] (5,12) circle (13);
  \draw[dashed] (5,0) -- (5,12);
  \draw[dashed] (0,12) -- (5,12);
  \fill (5,12) circle (\puntografico);
  \node[anchor=south east, inner sep=2pt, xshift=-2pt, yshift=2pt] at (5,12) {$T$};
  \node[anchor=north] at (5,0) {$5$};
  \node[left] at (0,12) {$12$};
\end{tikzpicture}
\end{center}

De acuerdo con la informaci\'on anterior, \textquestiondown en cu\'al de esos tel\'efonos se pueden realizar y recibir llamadas por medio de la se\~nal de esa antena?

\begin{enumerate}[label=\Alph*)]
  \item Solo en $R$.
  \item Solo en $S$.
  \item En $R$ y en $S$.
  \item En $R$, en $S$ y en $U$.
\end{enumerate}

\Needspace{0.88\textheight}
\item
En una finca se desea construir un estanque circular para peces. El centro del estanque ($C$) se ubicar\'a a $9$ m al oeste y $4$ m al sur de un \'arbol ($A$), el cual se considera como origen. Adem\'as, el borde del estanque debe quedar a $7$ m del centro.

La siguiente representaci\'on gr\'afica, en la que las unidades est\'an en metros, muestra la ubicaci\'on del \'arbol, del centro del estanque y de la circunferencia correspondiente al borde del estanque.

\begin{center}
\begin{tikzpicture}[scale=0.18]
  \path[use as bounding box] (-22,-16) rectangle (13,10);
  \draw[<->, thick] (-20,0) -- (11,0);
  \draw[<->, thick] (0,-14) -- (0,8);
  \node[anchor=west, inner sep=1pt] at (11.05,0) {$x$ (este)};
  \node[anchor=south, inner sep=1pt] at (0,8.55) {$y$ (norte)};
  \fill (0,0) circle (\puntografico) node[below right] {$A$};
  \draw[dashed] (-9,0) -- (-9,-4);
  \draw[dashed] (0,-4) -- (-9,-4);
  \draw[thick] (-9,-4) circle (7);
  \fill (-9,-4) circle (\puntografico) node[anchor=east, inner sep=2pt] {$C$};
  \node[anchor=north] at (-9,0) {$\llap{-}9$};
  \node[right] at (0,-4) {$-4$};
\end{tikzpicture}
\end{center}

\textquestiondown Cu\'al de las siguientes representaciones algebraicas, en la que las unidades est\'an en metros, corresponde a la circunferencia que representa el borde de ese estanque?

\begin{enumerate}[label=\Alph*)]
  \opcion{A}{(x-9)^2+(y-4)^2=49}
  \opcion{B}{(x+9)^2+(y+4)^2=49}
  \opcion{C}{(x+9)^2+(y+4)^2=7}
  \opcion{D}{(x-4)^2+(y-9)^2=49}
\end{enumerate}

\Needspace{0.68\textheight}
\item
En un plano cartesiano, la circunferencia $C$ tiene centro en $A(6,4)$ y radio $5$. A esa circunferencia se le aplica una traslaci\'on de $8$ unidades hacia la derecha y $3$ unidades hacia abajo, obteniendo la circunferencia $C'$.

\textquestiondown Cu\'al de las siguientes representaciones algebraicas corresponde a la circunferencia $C'$?

\begin{enumerate}[label=\Alph*)]
  \opcion{A}{(x-14)^2+(y-1)^2=25}
  \opcion{B}{(x+14)^2+(y+1)^2=25}
  \opcion{C}{(x-14)^2+(y-7)^2=25}
  \opcion{D}{(x-6)^2+(y-4)^2=25}
\end{enumerate}

\Needspace{0.90\textheight}
\item
En un centro recreativo, una piscina circular est\'a representada en un plano por la circunferencia $C$, con centro en $A(4,3)$ y radio $4$. Debido a una remodelaci\'on del terreno, la piscina se trasladar\'a $12$ unidades hacia la derecha y $5$ unidades hacia arriba, obteniendo la nueva ubicaci\'on representada por la circunferencia $C'$.

\textquestiondown Cu\'al de las siguientes representaciones gr\'aficas corresponde a esa traslaci\'on?

\vspace{2\baselineskip}

\begin{center}
\setlength{\tabcolsep}{0pt}
\renewcommand{\arraystretch}{1}
\begin{tabular}{@{}r@{\hspace{0.35cm}}c@{\hspace{0.55cm}}r@{\hspace{0.35cm}}c@{}}
\raisebox{1.45cm}{\textbf{A)}} & \graftras{4}{3}{4}{12}{8}{4} &
\raisebox{1.45cm}{\textbf{B)}} & \graftrasetiquetabaja{4}{3}{4}{16}{-2}{4} \\[0.75cm]
\raisebox{1.45cm}{\textbf{C)}} & \graftras{4}{3}{4}{16}{8}{4} &
\raisebox{1.45cm}{\textbf{D)}} & \graftras{4}{3}{4}{9}{15}{4}
\end{tabular}
\end{center}

\Needspace{0.90\textheight}
\item
En una mesa un carpintero traza una circunferencia a $6$ cm al este y $9$ cm al norte de uno de los extremos de la mesa, la cual representa la ubicaci\'on de un portavasos.

La siguiente representaci\'on gr\'afica, cuyas unidades est\'an en cm, muestra las ubicaciones, $A$ del centro de la circunferencia y $B$ extremo de la mesa, as\'i como la ubicaci\'on de la circunferencia correspondiente al alcance de la zona iluminada en el piso:

\begin{center}
\begin{tikzpicture}[scale=0.28]
  \path[use as bounding box] (-2,-2) rectangle (14,15);
  \draw[<->, thick] (-1,0) -- (12.5,0);
  \draw[<->, thick] (0,-1) -- (0,14);
  \node[anchor=west, inner sep=1pt] at (12.55,0) {$x$ (este)};
  \node[anchor=south, inner sep=1pt] at (0,14.15) {$y$ (norte)};
  \draw[dashed] (6,0) -- (6,9);
  \draw[dashed] (0,9) -- (6,9);
  \draw[dashed] (0,5) -- (6,5);
  \draw[thick] (6,9) circle (4);
  \fill (0,0) circle (\puntografico) node[above left] {$B$};
  \fill (6,9) circle (\puntografico) node[right] {$A$};
  \fill (6,5) circle (\puntografico);
  \node[anchor=north] at (6,0) {$6$};
  \node[left] at (0,5) {$5$};
  \node[left] at (0,9) {$9$};
\end{tikzpicture}
\end{center}

El carpintero traslada la circunferencia, desde su ubicaci\'on actual, a $30$ cm al norte y $40$ cm al oeste para luego trazar otra circunferencia que representar\'a otro portavasos.

De acuerdo con la informaci\'on anterior, \textquestiondown Cu\'al es la representaci\'on algebraica, cuyas unidades est\'an en cent\'imetros, de la circunferencia correspondiente a la circunferencia que corresponde al segundo portavasos?

\begin{enumerate}[label=\Alph*)]
  \opcion{A}{(x-46)^2+(y-39)^2=16}
  \opcion{B}{(x-39)^2+(y-46)^2=16}
  \opcion{C}{(x+24)^2+(y+31)^2=8}
  \opcion{D}{(x+34)^2+(y-39)^2=16}
\end{enumerate}

\Needspace{0.92\textheight}
\item
En las instalaciones de un centro recreativo hay una fuente circular de $16$ m de diametro. Si se crea un sendero rectilineo que pasa por los puntos $(20,4)$ y $(20,14)$.

La siguiente representacion grafica, cuyas unidades estan en metros, muestra el origen $O$ y $C$ del centro de la fuente:

\begin{center}
\begin{tikzpicture}[scale=0.34]
  \path[use as bounding box] (-2,-1.5) rectangle (22,19);
  \draw[<->, thick] (-1,0) -- (21,0);
  \draw[<->, thick] (0,-1) -- (0,18);
  \node[anchor=west, inner sep=1pt] at (21.05,0) {$x$ (este)};
  \node[anchor=south, inner sep=1pt] at (0,18.55) {$y$ (norte)};
  \fill (0,0) circle (\puntografico) node[below left] {$O$};
  \draw[dashed] (10,0) -- (10,9);
  \draw[dashed] (0,9) -- (10,9);
  \draw[thick] (10,9) circle (8);
  \fill (10,9) circle (\puntografico) node[anchor=south east, inner sep=2pt] {$C$};
  \node[anchor=north] at (10,0) {$10$};
  \node[left] at (0,9) {$9$};
\end{tikzpicture}
\end{center}

De acuerdo con la informacion anterior, la recta que representa ese sendero, con respecto a la circunferencia que representa el borde de esa fuente, es

\begin{enumerate}[label=\Alph*)]
  \item exterior.
  \item secante.
  \item tangente.
  \item Ninguna de las anteriores.
\end{enumerate}

\Needspace{0.96\textheight}
\item
Un juego consiste en lanzar una ficha desde cierta distancia hacia un objetivo circular colocado en el piso de un salon de actividades. La medida del radio de ese objetivo es $5$ m.

La siguiente representacion grafica, cuyas unidades estan en metros, muestra la circunferencia que representa el borde del objetivo:

\begin{center}
\begin{tikzpicture}[scale=0.38]
  \path[use as bounding box] (-2,-1.5) rectangle (17,16);
  \draw[<->, thick] (-1,0) -- (16,0);
  \draw[<->, thick] (0,-1) -- (0,15);
  \node[anchor=west, inner sep=1pt] at (16.05,0) {$x$ (este)};
  \node[anchor=south, inner sep=1pt] at (0,15.55) {$y$ (norte)};
  \draw[dashed] (9,0) -- (9,13);
  \draw[dashed] (0,8) -- (9,8);
  \draw[dashed] (0,13) -- (9,13);
  \draw[thick] (9,8) circle (5);
  \fill (9,8) circle (\puntografico);
  \fill (9,13) circle (\puntografico);
  \node[anchor=north] at (9,0) {$9$};
  \node[left] at (0,8) {$8$};
  \node[left] at (0,13) {$13$};
\end{tikzpicture}
\end{center}

Ademas, durante el juego:

\begin{itemize}
  \item Mariela lanzo una ficha y esta se ubico en el punto correspondiente a $(12,12)$.
  \item Andres lanzo una ficha y esta se ubico en el punto correspondiente a $(4,8)$.
  \item Valeria lanzo una ficha y esta se ubico en el punto correspondiente a $(15,9)$.
  \item Esteban lanzo una ficha y esta se ubico en el punto correspondiente a $(8,4)$.
\end{itemize}

De acuerdo con la informacion anterior, \textquestiondown cual de esas personas lanzo la ficha que se ubico en el exterior de ese objetivo?

\begin{enumerate}[label=\Alph*)]
  \item Mariela.
  \item Andres.
  \item Valeria.
  \item Esteban.
\end{enumerate}

\Needspace{0.96\textheight}
\item
Una rotonda tiene forma circular y la medida de su diametro es $24$ m. Las rectas $h$ y $j$ representan senderos que pasan por el centro de la rotonda. Las rectas $k$ y $n$ representan otros senderos del lugar ademas los senderos $j$ y $n$ son paralelos. La siguiente representacion grafica, en la que las unidades estan en metros, muestra la rotonda y esos senderos:

\begin{center}
\begin{tikzpicture}[scale=0.27]
  \path[use as bounding box] (-10,-7) rectangle (38,33);
  \draw[<->, thick] (-8,0) -- (36,0);
  \draw[<->, thick] (0,-5) -- (0,31);
  \node[anchor=west, inner sep=1pt] at (36.05,0) {$x$};
  \node[anchor=south, inner sep=1pt] at (0,31.05) {$y$};
  \draw[thick] (18,14) circle (12);
  \draw[very thick, <->] (18,-4) -- (18,32);
  \draw[very thick, <->] (2,14) -- (34,14);
  \draw[very thick, <->] (1,4) -- (35,4);
  \draw[very thick, <->] (6,30) -- (30,-2);
  \draw[thick] (18,4) -- (16.4,4) -- (16.4,5.6) -- (18,5.6);
  \draw[dashed] (0,14) -- (18,14);
  \draw[dashed] (18,0) -- (18,14);
  \draw[dashed] (0,4) -- (18,4);
  \fill (18,14) circle (\puntografico);
  \node[anchor=north] at (18,0) {$18$};
  \node[left] at (0,14) {$14$};
  \node[left] at (0,4) {$4$};
  \node[anchor=south east] at (18,32) {$h$};
  \node[anchor=south west] at (34,14) {$j$};
  \node[anchor=south east] at (8,29) {$k$};
  \node[anchor=south west] at (34,4) {$n$};
\end{tikzpicture}
\end{center}

De acuerdo con la informacion anterior, \textquestiondown cuales de esas rectas, que representan senderos, son perpendiculares entre si?

\begin{enumerate}[label=\Alph*)]
  \item $j$ y $k$.
  \item $h$ y $j$.
  \item $k$ y $n$.
  \item $j$ y $n$.
\end{enumerate}

\Needspace{0.92\textheight}
\item
Una antena de telecomunicaciones emite una senal que permite a los telefonos celulares, que se encuentran a $6$ km o menos alrededor de esa antena, realizar y recibir llamadas. A continuacion, se muestra la representacion grafica, en la que las unidades estan en kilometros, de la circunferencia que corresponde al alcance maximo de la senal que emite esa antena:

\begin{center}
\begin{tikzpicture}[scale=0.38]
  \path[use as bounding box] (-2,-1.5) rectangle (17,16);
  \draw[<->, thick] (-1,0) -- (16,0);
  \draw[<->, thick] (0,-1) -- (0,15);
  \node[anchor=west, inner sep=1pt] at (16.05,0) {$x$};
  \node[anchor=south, inner sep=1pt] at (0,15.05) {$y$};
  \draw[thick] (7,8) circle (6);
  \draw[dashed] (7,0) -- (7,14);
  \draw[dashed] (0,8) -- (13,8);
  \draw[dashed] (0,14) -- (7,14);
  \draw[dashed] (13,0) -- (13,8);
  \fill (7,8) circle (\puntografico);
  \fill (7,14) circle (\puntografico);
  \fill (13,8) circle (\puntografico);
  \node[anchor=north] at (7,0) {$7$};
  \node[anchor=north] at (13,0) {$13$};
  \node[left] at (0,8) {$8$};
  \node[left] at (0,14) {$14$};
\end{tikzpicture}
\end{center}

De acuerdo con la informacion anterior, si las ubicaciones que tienen los telefonos celulares $W$, $K$, $L$ y $M$ corresponden, respectivamente, a los puntos $(10,12)$, $(14,8)$, $(2,2)$ y $(7,15)$, entonces \textquestiondown en cual de esos telefonos se pueden realizar y recibir llamadas por medio de la senal de esa antena?

\begin{enumerate}[label=\Alph*)]
  \item $W$
  \item $K$
  \item $L$
  \item $M$
\end{enumerate}

\Needspace{0.92\textheight}
\item
En un centro educativo se destinara una zona con forma poligonal para construir un jardin. La siguiente representacion grafica, en la que las unidades estan en metros, muestra la forma de esa zona:

\begin{center}
\begin{tikzpicture}[scale=0.48]
  \path[use as bounding box] (-1.2,-1.2) rectangle (12,9.4);
  \draw[<->, thick] (-0.5,0) -- (11.5,0);
  \draw[<->, thick] (0,-0.5) -- (0,9);
  \node[anchor=west, inner sep=1pt] at (11.55,0) {$x$};
  \node[anchor=south, inner sep=1pt] at (0,9.05) {$y$};
  \fill[gray!18] (2,1) -- (10,1) -- (10,5) -- (6,8) -- (2,5) -- cycle;
  \draw[very thick] (2,1) -- (10,1) -- (10,5) -- (6,8) -- (2,5) -- cycle;
  \foreach \x/\lab in {2/2,10/10} {
    \draw[dashed] (\x,0) -- (\x,1);
    \node[anchor=north] at (\x,0) {$\lab$};
  }
  \foreach \y/\lab in {1/1,5/5} {
    \draw[dashed] (0,\y) -- (2,\y);
    \node[left] at (0,\y) {$\lab$};
  }
  \draw[dashed] (0,8) -- (6,8);
  \node[left] at (0,8) {$8$};
  \foreach \p/\pos/\name in {(2,1)/below left/A,(10,1)/below right/B,(10,5)/right/C,(6,8)/above/D,(2,5)/left/E} {
    \fill \p circle (\puntograficoitemunoafirmacioncuatro) node[\pos] {$\name$};
  }
\end{tikzpicture}
\end{center}

De acuerdo con la informacion anterior, \textquestiondown cuantos metros cuadrados mide el area de la zona destinada para ese jardin?

\begin{enumerate}[label=\Alph*)]
  \item $32$
  \item $40$
  \item $44$
  \item $56$
\end{enumerate}

\Needspace{0.70\textheight}
\item
Una empresa fabrica senales decorativas con forma de hexagono regular. Para colocar un refuerzo, se mide la distancia entre dos vertices opuestos de una de esas senales y se obtiene $18$ cm.

\begin{center}
\begin{tikzpicture}[scale=0.75]
  \path[use as bounding box] (-3.7,-3.2) rectangle (3.7,3.4);
  \draw[very thick] (0:3) -- (60:3) -- (120:3) -- (180:3) -- (240:3) -- (300:3) -- cycle;
  \draw[dashed, thick] (180:3) -- (0:3);
  \fill (0,0) circle (\puntograficoafirmacioncuatro);
  \node[below] at (0,0) {$O$};
  \node[above] at (0,0.35) {$18$ cm};
  \fill (180:3) circle (\puntograficoafirmacioncuatro);
  \fill (0:3) circle (\puntograficoafirmacioncuatro);
\end{tikzpicture}
\end{center}

Si se coloca una cinta alrededor de todo el borde de esa senal, \textquestiondown cuantos centimetros de cinta se necesitan?

\begin{enumerate}[label=\Alph*)]
  \item $27$
  \item $36$
  \item $54$
  \item $108$
\end{enumerate}

\Needspace{0.92\textheight}
\item
Un artesano diseno una pieza decorativa para la entrada de un salon comunal. La siguiente representacion grafica muestra la forma de esa pieza sobre una cuadricula. La medida del lado de cada cuadrado de la cuadricula es $1$ m.

\begin{center}
\begin{tikzpicture}[scale=0.58]
  \path[use as bounding box] (-0.8,-0.8) rectangle (9.3,8.4);
  \fill[gray!18] (2,1) -- (6,1) -- (6,5) arc[start angle=0,end angle=180,radius=2] -- cycle;
  \draw[step=1, black!70, dotted, line width=0.8pt] (0,0) grid (8,8);
  \draw[<->, thick] (-0.35,0) -- (8.4,0);
  \draw[<->, thick] (0,-0.35) -- (0,8.2);
  \node[anchor=west, inner sep=1pt] at (8.45,0) {$x$};
  \node[anchor=south, inner sep=1pt] at (0,8.25) {$y$};
  \draw[very thick] (2,1) -- (6,1) -- (6,5) arc[start angle=0,end angle=180,radius=2] -- cycle;
\end{tikzpicture}
\end{center}

De acuerdo con la informacion anterior, la superficie de esa pieza decorativa, en metros cuadrados, es mayor que

\begin{enumerate}[label=\Alph*)]
  \item $22$ y menor que $23$.
  \item $18$ y menor que $20$.
  \item $16$ y menor que $18$.
  \item $24$ y menor que $26$.
\end{enumerate}

\Needspace{0.88\textheight}
\item
Para una actividad escolar se arma un mural con cinco piezas cuadradas congruentes. Cada pieza tiene $4$ dm de lado y se colocan como se muestra en la siguiente figura:

\begin{center}
\begin{tikzpicture}[scale=0.75]
  \path[use as bounding box] (-0.5,-0.5) rectangle (4.5,4.5);
  \fill[gray!18] (1,0) rectangle (2,1);
  \fill[gray!18] (1,1) rectangle (2,2);
  \fill[gray!18] (0,1) rectangle (1,2);
  \fill[gray!18] (2,1) rectangle (3,2);
  \fill[gray!18] (1,2) rectangle (2,3);
  \draw[very thick] (1,0) -- (2,0) -- (2,1) -- (3,1) -- (3,2) -- (2,2) -- (2,3) -- (1,3) -- (1,2) -- (0,2) -- (0,1) -- (1,1) -- cycle;
  \draw[thin] (1,1) -- (2,1);
  \draw[thin] (1,2) -- (2,2);
  \draw[thin] (1,1) -- (1,2);
  \draw[thin] (2,1) -- (2,2);
  \draw[<->] (3.25,1) -- (3.25,2);
  \node[right] at (3.25,1.5) {$4$ dm};
\end{tikzpicture}
\end{center}

Si se coloca una cinta alrededor del borde exterior del mural, \textquestiondown cuantos decimetros de cinta se necesitan?

\begin{enumerate}[label=\Alph*)]
  \item $32$
  \item $40$
  \item $48$
  \item $80$
\end{enumerate}

\Needspace{0.42\textheight}
\item
En un taller se elaboran recuerdos a partir de esferas de acrilico. A cada esfera se le realiza un corte plano que pasa por el centro de la esfera.

De acuerdo con la informacion anterior, la seccion plana obtenida en cada recuerdo, producto del corte realizado a esa esfera, corresponde a una

\begin{enumerate}[label=\Alph*)]
  \item circunferencia.
  \item hiperbola.
  \item parabola.
  \item recta.
\end{enumerate}

\Needspace{0.48\textheight}
\item
Considere la siguiente informacion:

Una persona artesana trabaja con un bloque de cera con forma de cilindro circular recto, cuyo radio de la base mide $8$ cm y la altura del cilindro mide $24$ cm. Para elaborar una pieza decorativa, realiza un corte plano al bloque, que pasa por el centro de las bases y es perpendicular a estas.

De acuerdo con la informacion anterior, \textquestiondown cual es el perimetro de la seccion plana obtenida producto del corte realizado a ese bloque?

\begin{enumerate}[label=\Alph*)]
  \item $32$ cm
  \item $64$ cm
  \item $80$ cm
  \item $192$ cm
\end{enumerate}

\Needspace{0.78\textheight}
\item
Maria debe hacer cinco sombreros de cumpleanos como proyecto para un curso de fiestas infantiles. Ella confecciono los sombreros con un diametro de $14$ cm y una altura de $25$ cm. Si quiere pegar cinta formando una circunferencia paralela a la base y a $10$ cm de altura sobre la base del cono, como se muestra en la siguiente figura, \textquestiondown cuanta cinta debe comprar, aproximadamente, para los cinco sombreros?

\begin{center}
\begin{tikzpicture}[scale=0.62]
  \path[use as bounding box] (-4.4,-0.9) rectangle (4.7,7.3);
  \coordinate (V) at (0,6);
  \coordinate (L) at (-3,0);
  \coordinate (R) at (3,0);
  \coordinate (C) at (0,0);
  \coordinate (S) at (0,2.4);
  \coordinate (P) at (1.8,2.4);
  \draw[very thick] (V) -- (L);
  \draw[very thick] (V) -- (R);
  \draw[very thick] (0,0) ellipse (3 and 0.55);
  \draw[gray!40, very thick] (0,2.4) ellipse (1.8 and 0.33);
  \draw[dashed, thick] (0,0) -- (0,6);
\end{tikzpicture}
\end{center}

\begin{enumerate}[label=\Alph*)]
  \item $42$ cm
  \item $84$ cm
  \item $132$ cm
  \item $220$ cm
\end{enumerate}

\Needspace{0.98\textheight}
Considere la siguiente representaci\'on gr\'afica para responder los items 19 y 20.

En el plano de diseno de una institucion educativa, el poligono $ABCD$ representa la forma de una zona verde que sera reproducida en diferentes posiciones del patio.

\begin{center}
\begin{tikzpicture}[scale=0.48]
  \def\puntotransformacion{0.21}
  \path[use as bounding box] (-7.9,-1.2) rectangle (1.8,7.8);
  \draw[step=1, gray!45, very thin] (-7,0) grid (0,7);
  \draw[<->, thick] (-7.6,0) -- (1.3,0);
  \draw[<->, thick] (0,-0.7) -- (0,7.4);
  \node[anchor=west, inner sep=1pt] at (1.35,0) {$x$};
  \node[anchor=south, inner sep=1pt] at (0,7.45) {$y$};
  \fill[gray!15] (-2,6) -- (-4,1) -- (-6,4) -- (-5,5) -- cycle;
  \draw[densely dashed, gray!70, thin] (-2,0) -- (-2,6) -- (0,6);
  \draw[densely dashed, gray!70, thin] (-4,0) -- (-4,1) -- (0,1);
  \draw[densely dashed, gray!70, thin] (-6,0) -- (-6,4) -- (0,4);
  \draw[densely dashed, gray!70, thin] (-5,0) -- (-5,5) -- (0,5);
  \draw[very thick] (-2,6) -- (-4,1) -- (-6,4) -- (-5,5) -- cycle;
  \foreach \x/\lab in {-7/{\llap{-}7},-6/{\llap{-}6},-5/{\llap{-}5},-4/{\llap{-}4},-3/{\llap{-}3},-2/{\llap{-}2},-1/{\llap{-}1}}
    \node[anchor=north, inner sep=1pt] at (\x,0) {$\lab$};
  \foreach \y in {1,2,3,4,5,6,7}
    \node[anchor=east, inner sep=1pt] at (0,\y) {$\y$};
  \fill (-2,6) circle (\puntotransformacion) node[anchor=south west, inner sep=1pt, xshift=1pt, yshift=1pt] {\small A};
  \fill (-4,1) circle (\puntotransformacion) node[anchor=north east, inner sep=1pt, xshift=-1pt, yshift=-1pt] {\small B};
  \fill (-6,4) circle (\puntotransformacion) node[anchor=east, inner sep=1pt, xshift=-1pt] {\small C};
  \fill (-5,5) circle (\puntotransformacion) node[anchor=south east, inner sep=1pt, xshift=-1pt, yshift=1pt] {\small D};
\end{tikzpicture}
\end{center}

\item
Si al poligono $ABCD$ se le aplica una reflexion respecto a la recta dada por $y=-x$, entonces, \textquestiondown cual es la imagen del punto $A$?

\begin{enumerate}[label=\Alph*)]
  \item $(-2,-6)$
  \item $(6,-2)$
  \item $(-6,2)$
  \item $(2,6)$
\end{enumerate}

\item
En el mismo plano de diseno, al poligono $ABCD$ se le aplica una homotecia de razon $k=-2$ con centro en el origen y, luego, una traslacion de $3$ unidades hacia la derecha.

\textquestiondown Cual es la imagen del punto $C$ luego de aplicar ambas transformaciones?

\begin{enumerate}[label=\Alph*)]
  \item $(-15,8)$
  \item $(15,-8)$
  \item $(9,-8)$
  \item $(12,-5)$
\end{enumerate}

\Needspace{0.48\textheight}
\item
En una aplicacion de mapas, la ubicacion de una camara de seguridad se representa por el punto $P(3,-2)$ en un plano cartesiano. Para ajustar la orientacion del mapa en la pantalla, el sistema aplica una rotacion de $90^\circ$ en sentido antihorario alrededor del origen.

\textquestiondown Cuales son las coordenadas del punto que representa la nueva ubicacion de la camara en la pantalla?

\begin{enumerate}[label=\Alph*)]
  \item $(-2,-3)$
  \item $(2,3)$
  \item $(-3,2)$
  \item $(3,2)$
\end{enumerate}

\Needspace{0.98\textheight}
Considere la siguiente representacion grafica para responder los items 22 y 23.

En el siguiente plano cartesiano se muestra el poligono $MNPQR$.

\begin{center}
\begin{tikzpicture}[scale=0.55]
  \def\puntotransformacion{0.18}
  \path[use as bounding box] (-1.7,-6.3) rectangle (6.5,0.9);
  \draw[<->, thick] (-1.2,0) -- (5.8,0);
  \draw[<->, thick] (0,0.7) -- (0,-6.0);
  \node[anchor=west, inner sep=1pt] at (5.85,0) {$x$};
  \node[anchor=south, inner sep=1pt] at (0,0.72) {$y$};
  \foreach \x in {1,3,4}
    \draw[densely dotted, gray!60] (\x,0) -- (\x,-5);
  \foreach \y in {-1,-2,-4,-5}
    \draw[densely dotted, gray!60] (0,\y) -- (4,\y);
  \foreach \x in {1,3,4,5}
    \node[anchor=north, inner sep=2pt] at (\x,0) {$\x$};
  \foreach \y/\ylab in {-1/{$-1$},-2/{$-2$},-4/{$-4$},-5/{$-5$}}
    \node[anchor=east, inner sep=2pt] at (0,\y) {\ylab};
  \draw[thick] (1,-2) -- (4,-1) -- (4,-5) -- (1,-5) -- (3,-4) -- cycle;
  \fill (1,-2) circle (\puntotransformacion) node[anchor=south east] {$M$};
  \fill (4,-1) circle (\puntotransformacion) node[anchor=south west] {$N$};
  \fill (4,-5) circle (\puntotransformacion) node[anchor=north west] {$P$};
  \fill (1,-5) circle (\puntotransformacion) node[anchor=north east] {$Q$};
  \fill (3,-4) circle (\puntotransformacion) node[anchor=north west] {$R$};
\end{tikzpicture}
\end{center}

\item
Si al poligono $MNPQR$ se le aplica una reflexion respecto al eje $y$, \textquestiondown cual es la imagen del punto $N$?

\begin{enumerate}[label=\Alph*)]
  \item $(4, 1)$
  \item $(1, 4)$
  \item $(-1, -4)$
  \item $(-4, -1)$
\end{enumerate}

\item
Si al poligono $MNPQR$ se le aplica una homotecia de razon $k = -2$ con centro en el origen y, luego, una traslacion de $2$ unidades hacia la derecha, \textquestiondown cual es la imagen del punto $R$ luego de aplicar ambas transformaciones?

\begin{enumerate}[label=\Alph*)]
  \item $(-6, 6)$
  \item $(-4, 8)$
  \item $(-10, 8)$
  \item $(-6, 10)$
\end{enumerate}

\Needspace{0.90\textheight}
\item
En una estacion meteorologica se registra la relacion entre la variable $d$, que corresponde al dia de medicion, y la variable $r$, que corresponde a la cantidad de lluvia acumulada, en milimetros. Cada una de las siguientes tablas muestra una relacion obtenida durante tres semanas distintas.

\textbf{I.}

\begin{center}
\tablarel{$d$ & 1 & 2 & 3 & 2 & 5}{$r$ & 4 & 8 & 6 & 10 & 7}
\end{center}

\textbf{II.}

\begin{center}
\tablarel{$d$ & 0 & 1 & 2 & 3 & 4}{$r$ & 0 & 3 & 6 & 9 & 12}
\end{center}

\textbf{III.}

\begin{center}
\tablarel{$d$ & -2 & -1 & 0 & -2 & 1}{$r$ & 5 & 4 & 3 & 6 & 2}
\end{center}

\textbf{IV.}

\begin{center}
\tablarel{$d$ & 1 & 3 & 4 & 4 & 6}{$r$ & 2 & 5 & 7 & 9 & 12}
\end{center}

De acuerdo con la informacion anterior, \textquestiondown cual tabla muestra una relacion que corresponde a una funcion?

\begin{enumerate}[label=\Alph*)]
  \item I
  \item II
  \item III
  \item IV
\end{enumerate}

\Needspace{0.88\textheight}
\item
La siguiente representacion grafica corresponde al monto $C(x)$, en colones, que se debe pagar por el uso de un parqueo, en funcion del tiempo $x$, en horas, que permanece estacionado un vehiculo.

\begin{center}
\begin{tikzpicture}[xscale=0.55, yscale=0.65]
  \path[use as bounding box] (-1,-0.9) rectangle (11,6.2);
  \draw[->, thick] (-0.6,0) -- (10.5,0);
  \draw[->, thick] (0,-0.25) -- (0,5.9);
  \node[anchor=west, inner sep=1pt] at (10.55,0) {$x$};
  \node[anchor=south, inner sep=1pt] at (0,5.95) {$C(x)$};

  \foreach \x/\lab in {1/1,3/3,6/6,10/10}
    \node[anchor=north] at (\x,0) {$\lab$};
  \foreach \y/\lab in {0.8/800,1.5/{1500},3/{3000},5/{5000}}
    \node[anchor=east] at (0,\y) {$\lab$};

  \draw[densely dashed, gray!85, line width=0.55pt] (1,0) -- (1,0.8) -- (0,0.8);
  \draw[densely dashed, gray!85, line width=0.55pt] (3,0) -- (3,1.5) -- (0,1.5);
  \draw[densely dashed, gray!85, line width=0.55pt] (6,0) -- (6,3) -- (0,3);
  \draw[densely dashed, gray!85, line width=0.55pt] (10,0) -- (10,5) -- (0,5);

  \draw[thick] (0,0.8) -- (1,0.8);
  \draw[thick] (1,1.5) -- (3,1.5);
  \draw[thick] (3,3) -- (6,3);
  \draw[thick] (6,5) -- (10,5);

  \draw[densely dashed, black, line width=0.65pt] (1,0.8) -- (1,1.5);
  \draw[densely dashed, black, line width=0.65pt] (3,1.5) -- (3,3);
  \draw[densely dashed, black, line width=0.65pt] (6,3) -- (6,5);

  \puntoabierto{0}{0.8}
  \puntocerrado{1}{0.8}
  \puntoabierto{1}{1.5}
  \puntocerrado{3}{1.5}
  \puntoabierto{3}{3}
  \puntocerrado{6}{3}
  \puntoabierto{6}{5}
  \puntocerrado{10}{5}
\end{tikzpicture}
\end{center}

De acuerdo con la informacion anterior, si un vehiculo permanece estacionado durante $3$ horas, entonces, \textquestiondown cual monto se debe pagar por el uso del parqueo?

\begin{enumerate}[label=\Alph*)]
  \item 800 colones
  \item 1500 colones
  \item 3000 colones
  \item 5000 colones
\end{enumerate}

\Needspace{0.98\textheight}
Para responder los items 26 y 27, considere la siguiente informacion.

La siguiente representacion grafica corresponde a la funcion $n$, que determina la cantidad de botellas de agua disponibles, en cientos, en una bodega de emergencia, en funcion del tiempo $x$, en dias, desde el inicio de una jornada de distribucion, con $0 \leq x \leq 30$.

\begin{center}
\begin{tikzpicture}[xscale=0.20, yscale=0.16]
  \path[use as bounding box] (-6,-6) rectangle (36,33);
  \draw[->, thick] (-3,0) -- (34,0);
  \draw[->, thick] (0,-2) -- (0,31);
  \node[anchor=west, inner sep=1pt] at (34.05,0) {$x$};
  \node[anchor=south, inner sep=1pt] at (0,31.05) {$n(x)$};
  \node[rotate=90] at (-5,16) {Cantidad};
  \node at (17,-5.2) {Tiempo};

  \foreach \x/\lab in {5/5,10/10,15/15,20/20,25/25,30/30}
    \node[anchor=north] at (\x,0) {$\lab$};
  \foreach \y/\lab in {6/6,12/12,18/18,24/24,28/28}
    \node[anchor=east] at (0,\y) {$\lab$};

  \foreach \x/\y in {5/28,10/28,15/24,20/18,25/18,30/6}
    {
      \draw[densely dashed, gray!85, line width=0.55pt] (\x,0) -- (\x,\y);
      \draw[densely dashed, gray!85, line width=0.55pt] (0,\y) -- (\x,\y);
    }

  \draw[thick] (0,12) -- (5,28) -- (10,28) -- (15,24) -- (20,18) -- (25,18) -- (30,6);
  \foreach \x/\y in {0/12,5/28,10/28,15/24,20/18,25/18,30/6}
    \fill (\x,\y) circle (0.54);
\end{tikzpicture}
\end{center}

\item
\textquestiondown Cual fue la cantidad de botellas de agua disponibles en la bodega el dia $15$ desde el inicio de la jornada?

\begin{enumerate}[label=\Alph*)]
  \item 15 cientos
  \item 18 cientos
  \item 24 cientos
  \item 28 cientos
\end{enumerate}

\item
Durante todo ese tiempo, \textquestiondown cual fue la menor cantidad de botellas de agua disponibles en la bodega?

\begin{enumerate}[label=\Alph*)]
  \item 6 cientos
  \item 12 cientos
  \item 24 cientos
  \item 30 cientos
\end{enumerate}

\Needspace{0.62\textheight}
\item
En una aplicacion de entrenamiento, la funcion $f$ asigna a cada rutina su duracion recomendada, en minutos, y la funcion $g$ asigna a cada duracion la cantidad de puntos que obtiene una persona usuaria. Se sabe que
\[
f(x)=4x+10 \qquad \text{y} \qquad g(x)=2x-15.
\]

Si $x$ representa el nivel de una rutina, \textquestiondown cual es el valor de $(g\circ f)(5)$?

\begin{enumerate}[label=\Alph*)]
  \item 15
  \item 45
  \item 55
  \item 75
\end{enumerate}

\Needspace{0.95\textheight}
\item
La siguiente representacion grafica corresponde a la funcion $p$, que determina la presion $p(t)$, en kilopascales, registrada por un sensor industrial en funcion del tiempo $t$, en minutos, desde que se inicia una prueba, con $0 \leq t \leq 8$.

\begin{center}
\begin{tikzpicture}[scale=0.72]
  \path[use as bounding box] (-1.1,-0.9) rectangle (9.2,7.1);
  \draw[<->, thick] (-0.55,0) -- (8.8,0);
  \draw[<->, thick] (0,-0.45) -- (0,6.65);
  \node[anchor=west, inner sep=1pt] at (8.85,0) {$t$};
  \node[anchor=south east, inner sep=1pt] at (0,6.7) {$p(t)$};
  \node[anchor=north] at (8,0) {$8$};
  \node[anchor=east] at (0,2) {$2$};
  \node[anchor=east] at (0,6) {$6$};
  \draw[densely dashed, gray!85, line width=0.6pt] (8,0) -- (8,2) -- (0,2);
  \draw[thick] (0,6) -- (8,2);
  \fill (0,6) circle (0.18);
  \fill (8,2) circle (0.18);
  \node[anchor=south] at (3.2,4.55) {$p$};
\end{tikzpicture}
\end{center}

De acuerdo con la informacion anterior, \textquestiondown cual representacion grafica corresponde a $p^{-1}$?

\begin{center}
\setlength{\tabcolsep}{0pt}
\renewcommand{\arraystretch}{1}
\begin{tabular}{@{}r@{\hspace{0.45cm}}c@{\hspace{2.15cm}}r@{\hspace{0.45cm}}c@{}}
\raisebox{1.55cm}{\textbf{A)}} &
  \begin{tikzpicture}[scale=0.45, baseline=-0.5ex]
    \draw[<->, thick] (-0.4,0) -- (8.7,0);
    \draw[<->, thick] (0,-0.4) -- (0,6.7);
    \node[anchor=west, inner sep=1pt] at (8.75,0) {$x$};
    \node[anchor=south, inner sep=1pt] at (0,6.75) {$y$};
    \node[anchor=north] at (8,0) {$8$};
    \node[anchor=east] at (0,2) {$2$};
    \node[anchor=east] at (0,6) {$6$};
    \draw[densely dashed, gray!85, line width=0.6pt] (8,0) -- (8,2) -- (0,2);
    \draw[thick] (0,6) -- (8,2);
    \fill (0,6) circle (0.18);
    \fill (8,2) circle (0.18);
  \end{tikzpicture} &
\raisebox{1.85cm}{\textbf{B)}} &
  \begin{tikzpicture}[scale=0.45, baseline=-0.5ex]
    \draw[<->, thick] (-0.4,0) -- (8.7,0);
    \draw[<->, thick] (0,-0.4) -- (0,8.7);
    \node[anchor=west, inner sep=1pt] at (8.75,0) {$x$};
    \node[anchor=south, inner sep=1pt] at (0,8.75) {$y$};
    \node[anchor=north] at (2,0) {$2$};
    \node[anchor=north] at (6,0) {$6$};
    \node[anchor=east] at (0,8) {$8$};
    \draw[densely dashed, gray!85, line width=0.6pt] (6,0) -- (6,8) -- (0,8);
    \draw[thick] (2,0) -- (6,8);
    \fill (2,0) circle (0.18);
    \fill (6,8) circle (0.18);
  \end{tikzpicture} \\[1.05cm]
\raisebox{1.85cm}{\textbf{C)}} &
  \begin{tikzpicture}[scale=0.45, baseline=-0.5ex]
    \draw[<->, thick] (-0.4,0) -- (8.7,0);
    \draw[<->, thick] (0,-0.4) -- (0,8.7);
    \node[anchor=west, inner sep=1pt] at (8.75,0) {$x$};
    \node[anchor=south, inner sep=1pt] at (0,8.75) {$y$};
    \node[anchor=north] at (2,0) {$2$};
    \node[anchor=north] at (8,0) {$8$};
    \node[anchor=east] at (0,6) {$6$};
    \draw[densely dashed, gray!85, line width=0.6pt] (8,0) -- (8,6) -- (0,6);
    \draw[thick] (2,0) -- (8,6);
    \fill (2,0) circle (0.18);
    \fill (8,6) circle (0.18);
  \end{tikzpicture} &
\raisebox{1.85cm}{\textbf{D)}} &
  \begin{tikzpicture}[scale=0.45, baseline=-0.5ex]
    \draw[<->, thick] (-0.4,0) -- (8.7,0);
    \draw[<->, thick] (0,-0.4) -- (0,8.7);
    \node[anchor=west, inner sep=1pt] at (8.75,0) {$x$};
    \node[anchor=south, inner sep=1pt] at (0,8.75) {$y$};
    \node[anchor=north] at (2,0) {$2$};
    \node[anchor=north] at (6,0) {$6$};
    \node[anchor=east] at (0,8) {$8$};
    \draw[densely dashed, gray!85, line width=0.6pt] (2,0) -- (2,8) -- (0,8);
    \draw[thick] (2,8) -- (6,0);
    \fill (2,8) circle (0.18);
    \fill (6,0) circle (0.18);
  \end{tikzpicture}
\end{tabular}
\end{center}

\Needspace{0.56\textheight}
\item
Considere la siguiente informacion:

En un vivero se usa la funcion $a$, que determina la altura $a(s)$, en centimetros, de una planta de albahaca en funcion del tiempo $s$, en semanas, desde que se trasplanto, dada por
\[
a(s)=7s+18,
\]
donde $s$ representa el tiempo en semanas, con $2 \leq s \leq 9$.

De acuerdo con la informacion anterior, si se requiere determinar el criterio de la funcion inversa de $a$, entonces, \textquestiondown cual opcion corresponde a ese criterio?

\begin{enumerate}[label=\Alph*)]
  \item $s(a)=7a+18$
  \item $s(a)=\dfrac{a-18}{7}$
  \item $s(a)=\dfrac{a}{7}-18$
  \item $s(a)=\dfrac{18-a}{7}$
\end{enumerate}

\Needspace{0.56\textheight}
\item
Considere la siguiente informacion:

La funcion $q$, que determina la cantidad de litros de solucion nutritiva $q(r)$ que hay en un deposito, esta dada por
\[
q(r)=5r+12,
\]
donde $r$ representa el tiempo en horas desde que inicia el llenado del deposito, con $4 \leq r \leq 14$.

De acuerdo con la informacion anterior, si se define la funcion $r$, la cual corresponde a la funcion inversa de $q$, entonces, \textquestiondown cual es el ambito de $r$?

\begin{enumerate}[label=\Alph*)]
  \item $[4,14]$
  \item $[17,82]$
  \item $[32,82]$
  \item $[32,70]$
\end{enumerate}

\Needspace{0.58\textheight}
\item
Una plataforma de cursos registra una relacion entre cada codigo de curso y la cantidad de horas de tutoria asignadas por semana. Para que esa relacion tenga funcion inversa, \textquestiondown cual condicion debe cumplirse?

\begin{enumerate}[label=\Alph*)]
  \item Cada codigo de curso debe estar asociado con una unica cantidad de horas, aunque dos codigos diferentes tengan la misma cantidad.
  \item Cada cantidad de horas debe estar asociada con al menos dos codigos de curso.
  \item Cada codigo de curso debe estar asociado con una unica cantidad de horas y no debe haber dos codigos diferentes con la misma cantidad.
  \item Cada codigo de curso debe estar asociado con todas las cantidades de horas disponibles.
\end{enumerate}

\Needspace{0.58\textheight}
\item
Considere la siguiente informacion:

En un centro de acopio se estima la cantidad maxima $M$ de cajas pequenas que se pueden colocar en una zona de almacenamiento. Esa cantidad esta dada por
\[
M(x)=2\sqrt{x+9}+6,
\]
donde $x$ representa el area disponible adicional, en metros cuadrados, con $7 \leq x \leq 72$.

De acuerdo con la informacion anterior, \textquestiondown cual es la cantidad maxima de cajas pequenas que se pueden colocar si el area disponible adicional es de $16\text{ m}^2$?

\begin{enumerate}[label=\Alph*)]
  \item 11
  \item 14
  \item 16
  \item 31
\end{enumerate}

\Needspace{0.60\textheight}
\item
Considere la siguiente informacion:

En una planta de tratamiento se usa un sensor para estimar la cantidad $L$ de litros de agua filtrada por minuto. Esa cantidad esta dada por
\[
L(t)=3\sqrt{t+4}+2,
\]
donde $t$ representa el tiempo de funcionamiento del filtro, en minutos, con $0 \leq t \leq 60$.

De acuerdo con la informacion anterior, \textquestiondown cuantos minutos debe funcionar el filtro para que la cantidad estimada sea de $17$ litros por minuto?

\begin{enumerate}[label=\Alph*)]
  \item 9
  \item 21
  \item 25
  \item 29
\end{enumerate}

\Needspace{0.54\textheight}
\item
Considere la siguiente informacion:

En un laboratorio escolar se registra la altura $h(x)$, en centimetros, de una columna de espuma que se forma durante una reaccion. La variable $x$ representa el tiempo, en minutos, desde que inicia la observacion.

La siguiente representacion grafica muestra el comportamiento de esa altura:

\begin{center}
\begin{tikzpicture}[xscale=0.58, yscale=0.55]
  \path[use as bounding box] (-0.9,-0.8) rectangle (11.2,10.8);
  \draw[->, thick] (-0.5,0) -- (10.6,0) node[anchor=west] {$x$};
  \draw[->, thick] (0,-0.4) -- (0,10.2) node[anchor=south] {$h(x)$};

  \foreach \x/\lab in {1/1,2/2,5/5,10/10}
    \node[anchor=north] at (\x,0) {$\lab$};
  \foreach \y/\lab in {3/3,5/5,7/7,9/9}
    \node[anchor=east] at (0,\y) {$\lab$};

  \draw[densely dashed, gray!80] (1,0) -- (1,3) -- (0,3);
  \draw[densely dashed, gray!80] (2,0) -- (2,5) -- (0,5);
  \draw[densely dashed, gray!80] (5,0) -- (5,7) -- (0,7);
  \draw[densely dashed, gray!80] (10,0) -- (10,9) -- (0,9);

  \draw[thick, domain=1:10, samples=120, smooth, variable=\x]
    plot ({\x},{2*sqrt(\x-1)+3});

  \fill (1,3) circle (\puntograficoafirmacioncuatro);
  \fill (2,5) circle (\puntograficoafirmacioncuatro);
  \fill (5,7) circle (\puntograficoafirmacioncuatro);
  \fill (10,9) circle (\puntograficoafirmacioncuatro);
\end{tikzpicture}
\end{center}

De acuerdo con la informacion anterior, \textquestiondown cual representacion algebraica corresponde a la funcion $h$?

\begin{enumerate}[label=\Alph*)]
  \item $h(x)=2\sqrt{x-1}+3$
  \item $h(x)=2\sqrt{x+1}+3$
  \item $h(x)=2\sqrt{x-1}-3$
  \item $h(x)=\sqrt{x-1}+3$
\end{enumerate}

\Needspace{0.45\textheight}
\item
Considere la siguiente informacion:

Un sismografo registra la magnitud $M(t)$ de los movimientos sismicos en una zona cercana a una region volcanica del pais. La magnitud \textquotedblleft$M(t)$\textquotedblright, en grados, esta dada por
\[
M(t) = 1 + \log_3(t),
\]
donde \textquotedblleft$t$\textquotedblright\ es el tiempo en horas transcurridas desde el inicio del monitoreo, con $1 \leq t < 9$.

De acuerdo con la informacion anterior, la magnitud de los movimientos sismicos cuando han transcurrido mas de $3$ horas desde el inicio del monitoreo es

\begin{enumerate}[label=\Alph*)]
  \item mayor que $3$ grados y menor que $4$ grados
  \item mayor que $2$ grados y menor que $3$ grados
  \item mayor que $1$ grado y menor que $2$ grados
  \item mayor que $0$ grados y menor que $1$ grado
\end{enumerate}

\Needspace{0.42\textheight}
\item
Considere la siguiente informacion:

El tiempo $t$, en minutos, que necesita una impresora 3D para fabricar una pieza en funcion de la cantidad de capas $n$ de material utilizado esta dado por
\[
t = 2\log_4(n), \quad 1 \leq n \leq 64.
\]

De acuerdo con la informacion anterior, \textquestiondown cuantas capas de material se necesitan para que la impresora 3D tarde $6$ minutos en fabricar la pieza?

\begin{enumerate}[label=\Alph*)]
  \item $4$
  \item $16$
  \item $32$
  \item $64$
\end{enumerate}

\Needspace{0.56\textheight}
\item
La siguiente representacion grafica corresponde a la funcion $h$ que determina la cantidad de litros de agua $h(x)$ consumidos por hora en un sistema de riego, en funcion de la cantidad de aspersores $x$ activos, con $1 \leq x \leq 8$:

\begin{center}
\begin{tikzpicture}[xscale=0.65, yscale=0.65]
  \path[use as bounding box] (-1.2,-0.8) rectangle (9.5,13.5);
  \draw[->, thick] (-0.5,0) -- (9.0,0) node[anchor=west] {$x$};
  \draw[->, thick] (0,-0.4) -- (0,13.0) node[anchor=south] {$h(x)$};

  \foreach \x/\lab in {1/1,8/8}
    \node[anchor=north] at (\x,0) {$\lab$};
  \foreach \y/\lab in {5/5,12/12}
    \node[anchor=east] at (0,\y) {$\lab$};

  \draw[densely dashed, gray!80] (1,0) -- (1,5) -- (0,5);
  \draw[densely dashed, gray!80] (8,0) -- (8,12) -- (0,12);

  \draw[thick] (1,5) -- (8,12);

  \fill (1,5)  circle (\puntograficoafirmacioncuatro);
  \fill (8,12) circle (\puntograficoafirmacioncuatro);
\end{tikzpicture}
\end{center}

De acuerdo con la informacion anterior, \textquestiondown cual es el criterio de $h$?

\begin{enumerate}[label=\Alph*)]
  \item $h(x) = x$
  \item $h(x) = x + 4$
  \item $h(x) = 4x + 1$
  \item $h(x) = x - 4$
\end{enumerate}

\Needspace{0.40\textheight}
\item
Considere la siguiente informacion:

Una empresa de mensajeria cobra ₡$1\,500$ por cada paquete entregado, ademas de un costo fijo de ₡$5\,000$ por dia de operacion, sin importar la cantidad de paquetes entregados ese dia.

De acuerdo con la informacion anterior, el criterio que relaciona el ingreso total $I(x)$, en colones, de la empresa en funcion de la cantidad $x$ de paquetes entregados corresponde a

\begin{enumerate}[label=\Alph*)]
  \item $I(x) = 6\,500x$
  \item $I(x) = 5\,000x + 1\,500$
  \item $I(x) = 1\,500x + 5\,000$
  \item $I(x) = 1\,500x - 5\,000$
\end{enumerate}

\Needspace{0.44\textheight}
\item
Considere la siguiente informacion:

La funcion $v$ que determina la rapidez $v(t)$, en metros por segundo, de una lancha durante una competencia acuatica esta dada por
\[
v(t) = t^2 - 4t + 8,
\]
donde $t$ es el tiempo en minutos desde el inicio de la competencia, con $0 \leq t \leq 8$.

De acuerdo con la informacion anterior, si la rapidez de la lancha durante la competencia es de $20$ metros por segundo, \textquestiondown cual es el tiempo, en minutos, transcurrido desde el inicio de la competencia?

\begin{enumerate}[label=\Alph*)]
  \item $1$
  \item $2$
  \item $4$
  \item $6$
\end{enumerate}

\Needspace{0.42\textheight}
\item
Considere la siguiente informacion:

La funcion $g$ que determina la cantidad $g(t)$ de unidades vendidas por semana en un negocio esta dada por
\[
g(t) = -t^2 + 5t + 4,
\]
donde $t$ es el tiempo en semanas desde la apertura del negocio, con $1 \leq t \leq 5$.

De acuerdo con la informacion anterior, \textquestiondown cuantas unidades se vendieron en la semana $3$ desde la apertura del negocio?

\begin{enumerate}[label=\Alph*)]
  \item $4$
  \item $8$
  \item $10$
  \item $28$
\end{enumerate}

\Needspace{0.55\textheight}
\item
Considere la siguiente informacion:

En un taller artesanal se elaboran unicamente cajas de madera y estantes. Para su elaboracion, cada caja requiere 15 min de ensamblaje y 12 min de lijado. Asimismo, cada estante requiere 10 min de ensamblaje y 7 min de lijado.

De acuerdo con la informacion anterior, si para la elaboracion de estantes y cajas se utilizaron 345 min de ensamblaje y 264 min de lijado, entonces, \textquestiondown cuantas cajas de madera se fabricaron en total?

\begin{enumerate}[label=\Alph*)]
  \item 15
  \item 50
  \item 107
  \item 12
\end{enumerate}

\Needspace{0.48\textheight}
\item
Carlos y Mariana fueron a una cafeteria a comprar cafe molido y cafe en grano. Carlos compro 3~kg de cafe molido y 2~kg de cafe en grano, por los cuales pago \mbox{₡$24\,500$}. Mariana compro 2~kg de cafe molido y 4~kg de cafe en grano, por los cuales pago \mbox{₡$27\,000$}. Si cada kilogramo de cafe molido tenia el mismo precio y cada kilogramo de cafe en grano tenia el mismo precio, entonces, \textquestiondown cual fue el precio de cada kilogramo de cafe molido?

\begin{enumerate}[label=\Alph*)]
  \item \mbox{₡$7\,375$}
  \item \mbox{₡$5\,500$}
  \item \mbox{₡$4\,900$}
  \item \mbox{₡$4\,000$}
\end{enumerate}

\Needspace{0.92\textheight}
\item
Las empresas W y H cobran a sus respectivos clientes un monto por el alquiler de una bicicleta. En la siguiente representacion grafica se muestra el monto ``$y$'' en colones que cada empresa cobra a sus clientes, en funcion del tiempo ``$x$'' en horas, por el alquiler de la bicicleta.

\begin{center}
\begin{tikzpicture}[xscale=1.22, yscale=0.86]
  \path[use as bounding box] (-2.8,-1.3) rectangle (7.8,8.8);

  \draw[->, thick] (-0.3,0) -- (7.2,0) node[anchor=west, inner sep=1pt] {$X$};
  \draw[->, thick] (0,-0.5) -- (0,8.0) node[anchor=south, inner sep=1pt] {$Y$};

  \node[rotate=90, anchor=south, inner sep=6pt] at (-1.6,4.0) {Monto};
  \node[anchor=north, inner sep=6pt] at (3.5,-0.9) {Horas};

  \foreach \x in {2,4,6}
    \draw (\x,0.07) -- (\x,-0.07) node[anchor=north, inner sep=2pt, font=\small] {\x};

  \foreach \y/\lab in {1/1000, 2/2000, 3/3000, 3.5/3500, 5/5000, 6.5/6500, 7/7000}
    \draw (0.07,\y) -- (-0.07,\y) node[anchor=east, inner sep=2pt, font=\small] {\lab};

  \draw[densely dashed, gray!45] (2,0) -- (2,3.5);
  \draw[densely dashed, gray!45] (0,3.0) -- (2,3.0);
  \draw[densely dashed, gray!45] (0,3.5) -- (2,3.5);
  \draw[densely dashed, gray!45] (4,0) -- (4,5.0);
  \draw[densely dashed, gray!45] (0,5.0) -- (4,5.0);
  \draw[densely dashed, gray!45] (6,0) -- (6,7.0);
  \draw[densely dashed, gray!45] (0,6.5) -- (6,6.5);
  \draw[densely dashed, gray!45] (0,7.0) -- (6,7.0);

  % W: y = 1 + x (en miles de colones), intersecta H en x=4
  \draw[thick] (0,1) -- (6,7);
  \draw[thick, fill=white] (0,1) circle (0.10);
  \fill (2,3.0) circle (0.10);
  \fill (4,5.0) circle (0.10);
  \fill (6,7) circle (0.10);
  \node[anchor=south west, inner sep=2pt, font=\small\itshape] at (6.1,7.0) {$W$};

  % H: y = 2 + 0.75x (en miles de colones)
  \draw[thick] (0,2) -- (6,6.5);
  \draw[thick, fill=white] (0,2) circle (0.10);
  \fill (2,3.5) circle (0.10);
  \fill (6,6.5) circle (0.10);
  \node[anchor=north west, inner sep=2pt, font=\small\itshape] at (6.1,6.4) {$H$};
\end{tikzpicture}
\end{center}

De acuerdo con la informacion anterior, si un cliente de W alquilo una bicicleta por el mismo precio que la alquilo un cliente de H y ambos pagaron el mismo monto, entonces, \textquestiondown cuantas horas alquilo, cada uno de los clientes, la respectiva bicicleta?

\begin{enumerate}[label=\Alph*)]
  \item 2
  \item 4
  \item 8
  \item 6
\end{enumerate}

\Needspace{0.55\textheight}
\item
Considere la siguiente tabla referente a una funcion $h$.

\begin{center}
\begin{tabular}{|c|c|c|c|}
\hline
$x$    & 1  & 5  & 25 \\
\hline
$h(x)$ & 0 & $-1$ & $-2$ \\
\hline
\end{tabular}
\end{center}

De acuerdo con la informacion anterior, \textquestiondown cual es el criterio de $h$ que modela los datos de la tabla?

\begin{enumerate}[label=\Alph*)]
  \item $h(x) = 5^{x}$
  \item $h(x) = 5x^{2}$
  \item $h(x) = \log_{1/5}(x)$
  \item $h(x) = 5x$
\end{enumerate}

\Needspace{0.42\textheight}
\item
Considere la siguiente situacion:

Una fabrica de sombreros tiene un costo mensual fijo de \mbox{₡$400\,000$} (independientemente de la cantidad de sombreros que se produzca) y un costo adicional de \mbox{₡$250$} por cada sombrero confeccionado.

De acuerdo con la informacion anterior, el mejor modelo para representar la relacion entre el costo mensual de la fabrica y la cantidad de sombreros fabricados mensualmente, corresponde a una funcion

\begin{enumerate}[label=\Alph*)]
  \item lineal
  \item logaritmica
  \item exponencial
  \item cuadratica
\end{enumerate}

\Needspace{0.62\textheight}
\item
Considere la siguiente informacion:

En un laboratorio se estudia el crecimiento de una poblacion de bacterias. Su comportamiento se muestra en la siguiente tabla:

\begin{center}
\begin{tabular}{|p{5.5cm}|c|c|c|c|c|c|}
\hline
Cantidad de horas desde que se inicio el estudio. & 1 & 3 & 5 & 7 & 9 & 11 \\
\hline
Cantidad de bacterias & 4 & 12 & 28 & 52 & 84 & 124 \\
\hline
\end{tabular}
\end{center}

De acuerdo con la informacion anterior, si ``$x$'' representa la cantidad de horas desde que se inicio el estudio y ``$C(x)$'' corresponde a la cantidad de bacterias, entonces, \textquestiondown cual es el criterio que mejor modela la cantidad de bacterias, en funcion de la cantidad de horas, desde que se inicio el estudio?

\begin{enumerate}[label=\Alph*), itemsep=5pt]
  \item $C(x) = 2^{x} + 4$
  \item $C(x) = x^{2} + 3$
  \item $C(x) = \sqrt{x-3}$
  \item $C(x) = \log_{2}(x) + 4$
\end{enumerate}

\Needspace{0.82\textheight}
\item
La siguiente representacion grafica muestra la rapidez ``$R(x)$'', en metros por segundo, a la que corrio Miguel en una carrera, en funcion del tiempo ``$x$'' en minutos, desde el inicio de esa carrera, con $0 \leq x \leq 4$.

\begin{center}
\begin{tikzpicture}[xscale=1.5, yscale=1.5]
  \path[use as bounding box] (-0.9,-0.6) rectangle (5.5,3.0);
  \draw[->, thick] (-0.3,0) -- (5.0,0) node[anchor=west, inner sep=1pt] {$x$};
  \draw[->, thick] (0,-0.3) -- (0,2.7) node[anchor=south, inner sep=1pt] {$R(x)$};

  \draw (4,0.06) -- (4,-0.06) node[anchor=north, inner sep=2pt] {$4$};
  \draw[densely dashed, gray!70] (4,0) -- (4,2) -- (0,2);
  \node[anchor=east, inner sep=2pt] at (0,2) {$2$};

  \draw[thick, domain=0:4, samples=100, smooth, variable=\x]
    plot ({\x},{sqrt(\x)});

  \fill (0,0) circle (0.07);
  \fill (4,2) circle (0.07);
\end{tikzpicture}
\end{center}

De acuerdo con la informacion anterior, \textquestiondown cual de los siguientes criterios de funciones podria modelar la rapidez, en metros por segundo, a la que corrio Miguel en esa carrera, en funcion del tiempo en minutos, desde el inicio de esa carrera?

\begin{enumerate}[label=\Alph*)]
  \item $R(x) = 2^{x}$
  \item $R(x) = x^{2}$
  \item $R(x) = \sqrt{x}$
  \item $R(x) = \log_{2}(x)$
\end{enumerate}

\Needspace{0.82\textheight}
\item
La siguiente representacion grafica muestra la distribucion de los datos correspondiente a los tiempos, en horas, que tardo un grupo de ciclistas en recorrer una ruta de montana, durante un evento deportivo:

\begin{center}
\begin{tikzpicture}[xscale=1.8, yscale=9.5]
  \path[use as bounding box] (-0.5,-0.02) rectangle (5.8,0.44);
  \draw[->, thick] (-0.2,0) -- (5.5,0);

  \draw[thick, domain=0.005:5.2, samples=200, smooth, variable=\x]
    plot ({\x},{\x*exp(-\x)});

  \draw[densely dashed, gray!60] (2,0) -- (2,{2*exp(-2)});
  \node[anchor=north, inner sep=3pt, font=\small] at (2,0) {$2\,\mathrm{h}$};
\end{tikzpicture}
\end{center}

De acuerdo con la informacion anterior, si $2\,\mathrm{h}$ corresponde al promedio de esos datos, entonces la mediana de los tiempos que tardaron los ciclistas en recorrer esa ruta, es

\begin{enumerate}[label=\Alph*)]
  \item menor que $2\,\mathrm{h}$
  \item mayor que $2\,\mathrm{h}$
  \item igual que $2\,\mathrm{h}$
  \item Ninguna
\end{enumerate}

\Needspace{0.62\textheight}
\item
En la siguiente tabla se presentan algunas medidas de posicion referentes a las cantidades de dinero, en colones, ahorradas por las personas integrantes de un club social:

\begin{center}
\begin{tabular}{|l|r|}
\hline
\textbf{Medida de posicion} & \textbf{Valor} \\
\hline
Minimo         & 700\,000 \\
Maximo         & 1\,200\,000 \\
Moda           & 875\,000 \\
Mediana        & 900\,000 \\
Promedio       & 915\,000 \\
Primer Cuartil & 820\,000 \\
Tercer Cuartil & 1\,000\,000 \\
\hline
\end{tabular}
\end{center}

\textquestiondown Cual es la cantidad de dinero que con mayor frecuencia han ahorrado las personas integrantes de ese club?

\begin{enumerate}[label=\Alph*)]
  \item \mbox{₡$875\,000$}
  \item \mbox{₡$1\,000\,000$}
  \item \mbox{₡$1\,200\,000$}
  \item \mbox{₡$900\,000$}
\end{enumerate}

\Needspace{0.68\textheight}
\item
La siguiente tabla muestra el valor porcentual de cada uno de los tres componentes que se consideran al calificar proyectos de reciclaje de las personas participantes de un concurso. Asimismo, se muestra la calificacion obtenida por componente, en una escala de 1 a 100, del proyecto de reciclaje presentado por Cristian:

\begin{center}
\begin{tabular}{|l|c|p{4.5cm}|}
\hline
\textbf{Componente} & \textbf{Porcentaje} & \textbf{Calificacion del proyecto de Cristian} \\
\hline
Innovacion       & 45\% & \centering 89 \tabularnewline
Impacto ambiental & 35\% & \centering 85 \tabularnewline
Creatividad      & 20\% & \centering 73 \tabularnewline
\hline
\end{tabular}
\end{center}

Ademas, la nota final de ese proyecto corresponde a la media aritmetica ponderada de los tres componentes evaluados.

De acuerdo con la informacion anterior, \textquestiondown cual es la nota final del proyecto de reciclaje presentado por Cristian en ese concurso?

\begin{enumerate}[label=\Alph*)]
  \item $82{,}33$
  \item $84{,}40$
  \item $85{,}00$
  \item $84{,}00$
\end{enumerate}

\Needspace{0.65\textheight}
\item
Considere la siguiente informacion:

En la siguiente tabla se presentan algunas medidas de posicion referentes a los salarios mensuales, en colones, obtenidos por las personas trabajadoras de una empresa:

\begin{center}
\begin{tabular}{|l|r|}
\hline
\textbf{Medida de posicion} & \textbf{Valor} \\
\hline
Minimo         & 800\,000 \\
Maximo         & 1\,650\,000 \\
Moda           & 1\,100\,000 \\
Mediana        & 1\,250\,000 \\
Promedio       & 1\,310\,000 \\
Primer Cuartil & 950\,000 \\
Tercer Cuartil & 1\,500\,000 \\
\hline
\end{tabular}
\end{center}

Al menos el 75\% de los trabajadores de esa empresa obtuvo un salario mensual mayor o igual que

\begin{enumerate}[label=\Alph*)]
  \item \mbox{₡$950\,000$}
  \item \mbox{₡$1\,250\,000$}
  \item \mbox{₡$1\,500\,000$}
  \item \mbox{₡$1\,100\,000$}
\end{enumerate}

\Needspace{0.55\textheight}
\item
En el experimento de lanzar un dado legal y registrar el numero que sale en la cara superior, interesan dos eventos:

\begin{description}
  \item[A.] que el numero sea impar.
  \item[B.] que el numero sea un 4.
\end{description}

Con base en la informacion anterior, considere las siguientes proposiciones:

\begin{description}
  \item[I.] El evento A es el complemento del evento B.
  \item[II.] Los eventos A y B son mutuamente excluyentes.
\end{description}

\textquestiondown Cuales de ellas son verdaderas?

\begin{enumerate}[label=\Alph*)]
  \item Ambas
  \item Ninguna
  \item Solo I
  \item Solo II
\end{enumerate}

\Needspace{0.45\textheight}
\item
Considere el experimento de sacar dos bolas de una caja que contiene bolas blancas (B), azules (A) y rojas (R). Si las bolas poseen la misma forma y tamano, y el evento M es: que las bolas sean del mismo color, entonces con certeza el complemento $M^{C}$ de M es

\begin{enumerate}[label=\Alph*)]
  \item $\{AA,\; RR,\; BB\}$
  \item $\{BA,\; RB,\; BA,\; AB\}$
  \item $\{AA,\; BB,\; RR,\; BA,\; RA\}$
  \item $\{AB,\; AR,\; BA,\; BR,\; RA,\; RB\}$
\end{enumerate}

\Needspace{0.65\textheight}
Considere el siguiente contexto para responder los items 55 y 56.

\begin{center}
\textbf{Participantes de la carrera del Dia del Deporte, edicion 2014}

\vspace{4pt}
\begin{tabular}{|l|c|c|c|c|c|}
\hline
\textbf{Categoria} & \textbf{Juvenil} & \textbf{Mayor} & \textbf{Veterano A} & \textbf{Veterano B a D} & \textbf{Total} \\
\hline
Hombres & 18 & 96  & 42 & 28 & 184 \\
Mujeres & 6  & 56  & 32 & 14 & 108 \\
Total   & 24 & 152 & 74 & 42 & 292 \\
\hline
\end{tabular}
\end{center}

\item
\textquestiondown Cual es la probabilidad aproximada de elegir al azar a una participante mujer o de la categoria mayor?

\begin{enumerate}[label=\Alph*)]
  \item $0{,}37$
  \item $0{,}52$
  \item $0{,}70$
  \item $0{,}89$
\end{enumerate}

\item
\textquestiondown Cual es la probabilidad aproximada de elegir al azar a un participante hombre de la categoria veterano B a D o de la categoria veterano A?

\begin{enumerate}[label=\Alph*)]
  \item $0{,}24$
  \item $0{,}35$
  \item $0{,}40$
  \item $0{,}63$
\end{enumerate}

\Needspace{0.62\textheight}
Considere la siguiente informacion para responder los items 57 y 58.

En un colegio se registran las estaturas, en centimetros, de los estudiantes de cuatro equipos de baloncesto que compiten en un torneo institucional. El promedio de las estaturas y la desviacion estandar, por equipo, asi como la estatura de uno de los estudiantes de cada equipo, se presentan en la siguiente tabla:

\begin{center}
\begin{tabular}{|c|c|c|p{3.2cm}|}
\hline
\textbf{Equipo} & \textbf{Promedio} & \textbf{Desviacion estandar} & \centering\textbf{Estatura del estudiante} \tabularnewline
\hline
M & 174 & 6 & \centering 168 \tabularnewline
N & 170 & 5 & \centering 172 \tabularnewline
P & 175 & 8 & \centering 170 \tabularnewline
Q & 160 & 5 & \centering 164 \tabularnewline
\hline
\end{tabular}
\end{center}

\item
\textquestiondown En cual de los equipos se presenta una menor variabilidad relativa de los datos?

\begin{enumerate}[label=\Alph*)]
  \item Q
  \item P
  \item N
  \item M
\end{enumerate}

\item
\textquestiondown A cual equipo pertenece el estudiante que posee la mejor posicion relativa de su estatura, respecto a las estaturas de los estudiantes de su equipo?

\begin{enumerate}[label=\Alph*)]
  \item Q
  \item P
  \item N
  \item M
\end{enumerate}

\Needspace{0.92\textheight}
Considere la siguiente informacion para responder los items 59 y 60.

La siguiente tabla muestra informacion relacionada con los tiempos, en minutos, que tardaron diversos atletas de cuatro categorias al completar una carrera de 5~km:

\begin{center}
\begin{tabular}{|c|c|c|}
\hline
\textbf{Categoria} & \textbf{Media aritmetica} & \textbf{Desviacion estandar} \\
\hline
M & 21 & 7 \\
N & 27 & 9 \\
P & 31 & 10 \\
Q & 36 & 12 \\
\hline
\end{tabular}
\end{center}

Los datos obtenidos por cuatro atletas que participaron en la carrera se muestran a continuacion:

\begin{center}
\begin{tabular}{|l|c|c|}
\hline
\textbf{Nombre} & \textbf{Categoria} & \textbf{Tiempo en minutos} \\
\hline
Elena   & M & 39 \\
Julian  & N & 33 \\
Fabian  & P & 28 \\
Adriana & Q & 23 \\
\hline
\end{tabular}
\end{center}

\item
\textquestiondown Cual de los atletas obtuvo la mejor posicion relativa respecto a su categoria?

\begin{enumerate}[label=\Alph*)]
  \item Elena
  \item Julian
  \item Fabian
  \item Adriana
\end{enumerate}

\item
Considere las siguientes proposiciones:

\begin{description}
  \item[I.] Los tiempos de la categoria M presentan igual variabilidad relativa respecto a los tiempos de la categoria N.
  \item[II.] Los tiempos de la categoria P presentan menor variabilidad relativa respecto a los tiempos de la categoria Q.
\end{description}

De ellas, \textquestiondown cual o cuales son verdaderas?

\begin{enumerate}[label=\Alph*)]
  \item Ambas
  \item Ninguna
  \item Solo la I
  \item Solo la II
\end{enumerate}

\end{enumerate}

\end{document}
